% !PRINTED 210
\weberpage{210}
For the dihedral group of degree $n = 2m$, which we shall denote by $D_m$, we have, by the preceding section, two systems of $m$ conjugate two-fold poles each, and one system of two conjugate $m$-fold poles. The latter must therefore, by the theorem of \S.\,53, be the poles of a substitution. If we let them coincide with $0$ and $\infty$, we obtain the fundamental form of the second degree:
\[
f_1 = x_1 x_2.
\]
But this can be an invariant form of the group only if all the substitutions are contained in one of the two forms
\[
\begin{pmatrix} \alpha, & 0 \\ 0, & \delta \end{pmatrix}, \quad \begin{pmatrix} 0, & \beta \\ \gamma, & 0 \end{pmatrix}
\]
and here $\alpha:\delta$ and $-\beta:\gamma$ must be $m$th roots of unity. We therefore obtain the substitutions of the group, if $\varepsilon$ is a primitive $m$th root of unity:
\[
x, \varepsilon x, \varepsilon^2 x, \ldots, \varepsilon^{m-1} x
\]
\[
\frac{1}{x}, \frac{\varepsilon}{x}, \frac{\varepsilon^2}{x}, \cdots, \frac{\varepsilon^{m-1}}{x},
\]
i.\,e.
\weq{5}{D_m = \begin{pmatrix} \varepsilon^r, & 0 \\ 0, & 1 \end{pmatrix}, \begin{pmatrix} 0, & \varepsilon^r \\ 1, & 0 \end{pmatrix},}
or, represented with the determinant $+1$:
\weq{6}{D_m = \begin{pmatrix} \varepsilon^{1/2\,r}, & 0 \\ 0, & \varepsilon^{-1/2\,r} \end{pmatrix}, \begin{pmatrix} 0, & i\varepsilon^{1/2\,r} \\ i\varepsilon^{-1/2\,r}, & 0 \end{pmatrix}, \quad r = 0, 1, \ldots, m-1.}
One sees at once that these substitutions really form a group.

Besides $0$ and $\infty$ we have here also the poles $\pm \varepsilon^{1/2\,h}$ arising from the equations
\[
x = \frac{\varepsilon^h}{x}
\]
from which one further obtains the two fundamental forms of degree $m$:
\weq{7}{f_2 = x_1^m + x_2^m, \quad f_3 = x_1^m - x_2^m}

From this one sees that the dihedral groups and the cyclic groups are entirely different from one another, since their fundamental forms have different degrees.
